% DELTA23 root proof draft. All uses of E.58/E.65 are explicit supplied inputs.
\subsection{What a two--three Hilbert-symbol interpretation proves}

Write $\mathcal G=\widehat{GT}_0$, $\mathcal H=\ker\lambda$, and
$\mathcal H_5=\widehat{GT}\cap\mathcal H$. We retain the signed
frame coordinates of LM7\_PUSH E.65 and the CRT map of
DELTA\_ON\_GT A.6. Thus $a_2=\epsilon_2(W)$,
$a_3=\epsilon_3(T)$ and $\delta=a_2+a_3$ on $\mathcal H$.
The additive character $\Delta$ and the multiplicative Dirichlet
character $\chi_{12}$ have different codomains: the correctly typed
identity is $(-1)^{\delta(h)}=\chi_{12}(u(h))$.

\paragraph{D23.L1 (a well-defined Hilbert realization of the compressed frames).}
For $h\in\mathcal H$, choose any rational number $n$ which is a unit
at $2$ and $3$ and has residues
\[
n\equiv(-1)^{a_2(h)}\pmod4,\qquad
n\equiv(-1)^{a_3(h)}\pmod3.
\]
In particular one may choose the unique $n\in\{1,5,7,11\}$.
Define $L_2(h)=(n,3)_2$ and $L_3(h)=(n,3)_3$, using the quadratic
Hilbert symbols over $\mathbf Q_2$ and $\mathbf Q_3$. These are
independent of the indicated choice of $n$, are continuous characters,
and satisfy
\[
L_2=(-1)^{a_2},\qquad L_3=(-1)^{a_3},\qquad
L_2L_3=(-1)^\delta=\chi_{12}(u).
\tag{D23.L1.1}
\]
For the E.58 element $h_*$, in the supplied convention
$(a_3,a_2)=(1,0)$, one has $n=5$ and $(L_2,L_3)=(1,-1)$.
\begin{proof}
Use the standard local Hilbert formulas, with the precise locators
Voight, \emph{Quaternion algebras}, (12.4.10), 12.4.12(b), and
(12.4.14). For a $2$-adic unit $n$ they give
$(n,3)_2=(-1)^{(n-1)/2}$, where the exponent is read modulo two.
For a $3$-adic unit $n$ they give $(n,3)_3=(n/3)$.
These depend only on $n\bmod4$ and $n\bmod3$, respectively, and
give the first two identities. Their product gives the third by
the supplied definition of $\delta$. The four residues give
$\chi_{12}$ with values $1,-1,-1,1$, in the order $1,5,7,11$.
Continuity and multiplicativity follow from the already proved
frame characters, not from a purported multiplicative choice of
rational representatives. Finally E.58.23 and E.65.2--4 give
$T=-I$ at three and $W=I$, so the last evaluation follows.
\end{proof}

This is a positive answer for an explicitly defined pair of symbols
of the \emph{compressed} frame data. The definition has used the
two frame parities. It is not an independent arithmetic comparison
and does not prove that the two symbols have equal values.

\paragraph{D23.L2 (the proposed failure of global reciprocity is false).}
The rational pair $(n,3)$ in D23.L1 obeys Hilbert reciprocity for
every frame value. For the four positive representatives its complete
ramification table is
\[
\begin{array}{c|cc|c|c}
n&(n,3)_2&(n,3)_3&\text{other ramified finite primes}&(n,3)_\infty\\\hline
1&1&1&\varnothing&1\\
5&1&-1&5&1\\
7&-1&1&7&1\\
11&-1&-1&\varnothing&1
\end{array}
\tag{D23.L2.1}
\]
In particular E.58 does not violate a global product formula: its
extra factor is $(5,3)_5=-1$.
\begin{proof}
At an odd prime not dividing $3n$, both arguments are units and
the Hilbert symbol is one by Voight 12.4.12(a). At a prime dividing
one of $5,7,11$ the symbol is $(3/p)$. The square residues give
$(3/5)=(3/7)=-1$ and $(3/11)=1$ (since $5^2=3\pmod{11}$).
The two local values are D23.L1, and the real symbol is one
because $3>0$. These facts give the entire table. Each row has
an even number of minus signs, in agreement with Hilbert
reciprocity, Voight Proposition 14.2.1 (or Corollary 14.6.2).
\end{proof}

\paragraph{D23.L3 (the missing support hypothesis is exactly the old parity gate).}
For two bits $(b_2,b_3)$, the following are equivalent:
\begin{enumerate}
\item $b_2=b_3$;
\item there is $a\in\mathbf Q^\times$, a unit at $2$ and $3$, with
$(a,3)_2=(-1)^{b_2}$ and $(a,3)_3=(-1)^{b_3}$, and with
$(a,3)_v=1$ at every other place;
\item there is a rational quaternion algebra whose local invariants
at $2,3$ are $b_2/2,b_3/2$ and which is split at every other place.
\end{enumerate}
For the two equal-bit cases the algebras can be chosen as
$M_2(\mathbf Q)$ and the discriminant-six algebra $(2,3)$.
Thus the existence of such a two-place globalization for every
$h\in\mathcal H_5$ is \emph{equivalent} to A.84; it is not an
already proved arithmetic reason for A.84.
\begin{proof}
The implication from the third condition to the first is Hilbert
reciprocity (or the even-ramification part of Voight Main Theorem
14.6.1). The second condition implies the third by taking $(a,3)$.
For equal bits choose $a=1$ or $a=-1$. The local formulas give,
for $a=-1$, minus signs exactly at $2$ and $3$, and the real symbol
is again positive. This proves the first implication to the second.
Finally $(-1,3)$ is isomorphic to $(2,3)$. For an explicit change
of generators in $(-1,3)$, write $i^2=-1$, $j^2=3$, $ij=-ji$.
Then $i'=(1+j)i$ has square
$(1+j)(1-j)i^2=(1-3)(-1)=2$ and anticommutes with $j$.
Since $1+j$ is invertible, $i',j$ generate the original algebra.
This proves the asserted identification without any appeal merely
to coincident discriminants. Applying the equivalence to
$(b_2,b_3)=(a_2(h),a_3(h))$ gives the final assertion.
\end{proof}

\paragraph{D23.L4 (relative pro-three frames cannot recover the cubic bit).}
Let $U=F'F^3$ and let $\Gamma_3=B/O^3(U)$ be the specified
relative maximal pro-three quotient used in E.58. There is no
function of the induced automorphism of $\Gamma_3$ together with
the two retained pure frame values there which equals $a_3$ on
all of $\mathcal H$. In particular the cubic bit cannot be a
normalized local symbol depending only on that data.
\begin{proof}
The identity element has identity action and identity frames, with
$a_3=0$. The actual element supplied by E.58 has identity action
on every relative local factor. By E.58.20 its two frame values
in the factor at three are both one, while E.58.23 gives $a_3=1$.
The input data for the proposed function are equal and the
required outputs differ. This is a contradiction.
\end{proof}

The quotient $\Gamma_3$ in D23.L4 must not be confused with the
matrix target $\operatorname{SL}_2(\mathbf Z_3)$. The latter has
non-three-primary finite subgroups; reduction modulo three already
retains a quaternion group. D23.L4 refutes a precise pro-three
interpretation, not every possible arithmetic local comparison.

\paragraph{D23.L5 (full multipliers and the cyclotomic distinction).}
For $g\in\mathcal G$ choose $\sigma\in G_{\mathbf Q}$ with
$\chi(\sigma)=\lambda(g)$ and put $h=g\sigma^{-1}$.
The product $L_2(h)L_3(h)=(-1)^{\Delta(g)}$ is independent of this
choice, while the individual symbols need not be. Moreover neither
of the two individual bits is the restriction of a Dirichlet
character of the multiplier.
\begin{proof}
Choice independence of the product is E.65.5--6. The source
DELTA\_ON\_GT A.25 constructs $\sigma_0\in G_{\mathbf Q^{\rm ab}}$
which sends $r_0=(3/4)^{1/4}$ to $-r_0$, with
$a_2(\sigma_0)=a_3(\sigma_0)=1$ and $\lambda(\sigma_0)=1$.
For $g=1$, taking $\sigma=1$ gives the two symbols $(1,1)$,
whereas taking $\sigma=\sigma_0$ gives $(-1,-1)$. This proves
the individual ambiguity. A character of the multiplier is
trivial on every multiplier-one element, while each of these
frame bits takes value one at $\sigma_0$. This proves the last
assertion. The cited example concerns a nonabelian quartic-root
extension over $\mathbf Q$, quadratic only after adjoining
$\mathbf Q^{\rm ab}$; it is not the cyclotomic quadratic field
$\mathbf Q(\sqrt3)$.
\end{proof}
